\documentclass{zhvt-classic}

% 使用本地源流明体（GenRyuMin2 TW）按文件路径加载，不依赖系统字体名
% 版框（内框／界栏／书口侧栏／鱼尾）改用与版心印章同一印泥棕色：
%   印章图片 yshs.png 实色 = RGB(135,68,52) = HTML 874434
\usepackage{xcolor}
\usepackage{tikz}
% sealbrown 双注册：xcolor 供 TikZ/pgf，L3 供 cls 版框
\definecolor{sealbrown}{HTML}{874434}
\ExplSyntaxOn
\color_set:nnn { sealbrown } { HTML } { 874434 }
\ExplSyntaxOff

\zhvtset{
  grid color       = sealbrown,
  main font        = GenRyuMin2TW-SB,
  main font path   = C:/Users/Administrator/texmf/fonts/opentype/public/GenRyuMin2TW/,
jiazhu font      = GenRyuMin2TW-M,
  jiazhu font path = C:/Users/Administrator/texmf/fonts/opentype/public/GenRyuMin2TW/,
  % 页码框容纳 4 字：正常页「九十七」+前置空 1 格=4 行恰好；
  % 100+ 页「一百零一」共 5 行超出，末字「一/二/三」被框裁掉（2026-10-04 实测），故放宽到 5
  shukou page number chars = 5,
  % 书口书名框高从默认 6 字减到 4 字，把省出的 2 字高拨给 chapter 框，
  % 使 6 字书名+卷名放进 9 字高的 chapter 框后不再压线（2026-10-05）
  shukou title chars = 4,
}

\title{史记}[儒林列传]
\author{汉·司马迁}
\maker{闲叟漫记·重製}


% ===== 索隐述赞样式（2026-11-08 四次定稿 · 愚大人亲手改定参数）=====
% ⚠ 竖排前提：自定义字体家族必须带 RawFeature={vertical:+vert}；
%   simhei.ttf 加 +vhal 会整体加载失败静默回退，只写 +vert；simfang 同理。
\newCJKfontfamily{\shuzanfont}[Path=C:/Windows/Fonts/,RawFeature={vertical:+vert}]{simhei.ttf}
\newCJKfontfamily{\shuzanbody}[Path=C:/Windows/Fonts/,RawFeature={vertical:+vert}]{simfang.ttf}
% 题头：152x30pt 圆角 6pt 棕底白字；baseline (0,2) 使整框落入界行格子不压版框线；
% 白字节点 x 不补偿、y +8.2pt（像素级墨迹实测 dx/dy≈0，L/R 留白约 20pt）。
\newcommand{\shuzanhead}[1]{%
  \begin{tikzpicture}[baseline={(0,2)}]
    \fill[sealbrown, rounded corners=6pt]
      (-76pt,-15pt) rectangle (76pt,15pt);
    \node[text=white, inner sep=0pt, anchor=center] at (0pt,8.2pt)
      {\shuzanfont\fontsize{23}{23}\selectfont #1};
  \end{tikzpicture}%
}
% ★ 修正书口卷名位置：chapter 框起始位置改回 cls 默认值，避免全局影响。
%   只在每个卷首页通过 \juan 命令里的空命令微调该页位置（安全，不会跨页污染）。
\ExplSyntaxOn
\AtBeginDocument{
  \dim_set:Nn \l__zhvt_shukou_chapter_title_start_dim
    { \l__zhvt_shukou_title_height_dim + \l__zhvt_font_size_dim * 2 }
  \dim_set:Nn \l__zhvt_shukou_page_number_start_dim
    { \l__zhvt_shukou_chapter_title_start_dim + \l__zhvt_shukou_chapter_title_height_dim }
}
\ExplSyntaxOff

\begin{document}

% ★★ 版心（书口）排布覆写：把书名放到鱼尾之下（2026-10-04）
%
%   cls 默认版心自上而下：书名（title 静态框，在鱼尾之上）→ 鱼尾 → 篇名（chapter 动态框）→ 页码
%   本模板改为：          鱼尾（位置不动）→ 书名 → 篇名 → 页码分割线 → 页码
%
%   手法：把 cls 放在 title 静态框里的书名，改放进 chapter 动态框（鱼尾之下）；
%         title 静态框留空。书名取 \l__zhvt_title_str，即 \title{...} 的第一参数，
%         所以换书时只改 \title，不必在此重复写书名。
%
%   ⚠ 三条必须遵守的写法（2026-10-04 实测）：
%     ① 汉字必须用 \hbox{...} 包住。flowfram 会把框内容装进一个宽 14.4mm（= 一个 line_sep
%        = 书口栏宽）的纵向 minipage 并按该宽度折行，而整页靠 PDF /Rotate 90 显示、
%        框内文字本身不变向 —— 不包 \hbox 时「围炉夜话」4 字会被逐字折成 4 行，
%        在版面上就成了横排 4 字并横向溢出版框。
%     ② 竖排下**一个 \hbox 就是一整列**。书名与篇名要同处一列，必须放进同一个 \hbox；
%        分成两个 \hbox 时，后一个会被当成新的一列、甩到隔壁栏去
%        （实测「序」跑到书口栏外侧 14.4mm 的正文列上）。中间的 \quad 在框内是横向空白，
%        显示出来正好是书名与篇名之间的纵向间隔。
%     ③ ★ 内容盒要写成 **与框等宽的定宽盒** `\hbox to \l__zhvt_shukou_chapter_title_height_dim{…\hfil}`。
%        实测：chapter 框是把内容盒沿纵向**居中**摆放的，内容宽度一变，书名就跟着漂：
%        自然宽度的 \hbox（书名 38.4mm）落在 x=157.2mm，加篇名后（57.6mm）落在 x=152.4mm，
%        目录页（无篇名）与正文页的书名就差半格。做成等宽盒后，无论有没有篇名，
%        内容固定从框顶起排，实测两页均落在 x=142.8mm，不再漂。
\ExplSyntaxOn
% 把书口 chapter 框（书名+篇名）起始位置下移 1 字高，避免 4 字卷名顶到鱼尾
% ⚠ cls 的 \__zhvt_calculate_dimensions: 在 \AtBeginDocument 里执行，会重算该 dim；
%   必须把覆写挂到 begin document 之后，否则会被二次覆盖（2026-10-05 实测）。
% (旧方案已移到导言区，直接修改 chapter 框起始位置)
% 卷标题命令 #1=书口篇名(鱼尾) #2=正文卷名（含卷序，卷名内空格用 ~）
%
% ★★ 目录页两项修订（2026-10-05，愚大人指出）★★
%   ① 页码起点统一对齐：
%      原状 —— 目录项写成 \par\hspace{4em}卷名\hspace{4em}页码，页码起点 = 卷名实宽 + 4em。
%      卷名实宽随字数变：「卷一~梁惠王上」6 字 +~、「卷七~离娄上」5 字 +~，
%      差 1 个字位 → 卷七~卷十 4 条的页码整体比其余 10 条高 27.21pt（一个字），看着上下不齐。
%      改法 —— 卷名补全角空格（宽 = 1 字 = 27.215pt）补足到 **7 字符位**再写进目录，
%      14 条页码起点遂全部落在同一条竖线上（起点对齐；末字随页码字数自然参差）。
%      ⚠ 只对目录里的副本 \l_tmpc_tl 补，正文卷名仍用 \l_tmpa_tl 原样 —— 正文像素不受影响。
%   ② 卷名整体上移半字位：前导 \hspace 由 4em 改 3.5em。
%      实测原状目录块在版框内偏下（天头留白 139.6pt / 地脚留白 110.7pt），
%      上移 0.5em = 13.6pt 后上下基本均分，视觉居中。
\NewDocumentCommand { \juan } { mm } {
	\clearpage
	\tl_set:Nn \l__zhvt_chapter_title_tl { #1 }
	\tl_set:Nn \l_tmpa_tl { #2 }
	\tl_set_eq:NN \l_tmpc_tl \l_tmpa_tl
	\int_step_inline:nn { 7 - \str_count:n { #2 } }
	  { \tl_put_right:Nn \l_tmpc_tl { 　 } }
	\tl_set:Nx \l_tmpb_tl { \zhnum { page } }
	\tl_put_left:Nn \l_tmpb_tl { \hspace { 4em } }
	\tl_put_left:Nn \l_tmpb_tl { \l_tmpc_tl }
	\tl_put_left:Nn \l_tmpb_tl { \par \hspace { 3.5em } }
	\addtocontents { toc } { \l_tmpb_tl }
	\skip_horizontal:n { \l__zhvt_font_size_dim * 4 }
	\l_tmpa_tl
	\par
	\thispagestyle { empty }
}
\cs_set_protected:Npn \__zhvt_set_shukou_content: {
	% ★ 鱼尾上方留白（title 静态框）嵌入闲叟钤印：
	%   整页 /Rotate 90，故此处把图用 \rotatebox{90} 转正为竖章；
	%   尺寸卡在版心条幅内（书内约 7mm 宽 × 33.8mm 高，留余量不压鱼尾）。
	\setstaticcontents * { title } {
		\vbox_to_ht:nn { \l__zhvt_baseline_skip_dim } {
			\vskip 14.2pt
			\hbox_to_wd:nn { \l__zhvt_shukou_title_height_dim } {
				\hfil
				\rotatebox [ origin = c ] { 90 } { \includegraphics [ width = 7mm , height = 33.8mm ] { yshs.png } }
				\hfil
			}
		}
	}
% ★ 版心（书口）统一改法（2026-11-16 愚大人定）：18pt＋框内居中＋书口文字同版框色
%   ① 原 27.2pt 时 7 字长篇名（如「廉颇蔺相如列传」）末字溢出 chapter 框、压住页码；
%      改 \fontsize{18pt}{18pt}，前后 \hfil 令字串在框内居中（不随篇名字数漂移）。
%   ② \vbox{\vskip 8.2pt} 定量下移，抵消缩字号后栏宽方向 8.2pt 上偏（像素实测 y 中心 543.3 = 原版）。
%      ⚠ \raisebox / \vspace* 在 flowfram 框内无效；换字号须同步换补偿：16pt→10.1、19pt→7.4。
%   ③ chapter 与 page~number 均置 \color{sealbrown}（HTML 874434），与版框／界行同色；正文仍黑。
%   ⚠ 页码字位、钤印、鱼尾、界行、正文一律未动；五部合订本同此口径。
	\setdynamiccontents * { chapter }
	{ \color{sealbrown} \vbox { \vskip 8.2pt \hbox to \l__zhvt_shukou_chapter_title_height_dim
		{ \hfil \fontsize{18pt}{18pt}\selectfont \l__zhvt_title_str \quad \l__zhvt_chapter_title_tl \hfil } } }
	\setdynamiccontents * { page~number } { \color{sealbrown} \quad \zhnum { page } }
}
\__zhvt_set_shukou_content:
\ExplSyntaxOff

\maketitle

% 单篇，无多卷目录，不渲染目录页

\juan{儒林列传}{卷一百二十一~儒林列传第六十一}
\hspace*{2em}太史公曰。余读功令。至於广厉学官之路。未尝不废书而叹也。曰。嗟乎。夫周室衰而关雎作。幽厉微而礼乐坏。诸侯恣行。政由彊国。故孔子闵王路废而邪道兴。於是论次诗书。修起礼乐。適齐闻韶。三月不知肉味。自卫返鲁。然後乐正。雅颂各得其所。世以混浊莫能用。是以仲尼干七十馀君无所遇。曰﹁苟有用我者。期月而已矣﹂。西狩获麟。曰﹁吾道穷矣﹂。故因史记作春秋。以当王法。其辞微而指博。後世学者多录焉。\\
\hspace*{2em}自孔子卒後。七十子之徒散游诸侯。大者为师傅卿相。小者友教士大夫。或隐而不见。故子路居卫。子张居陈。澹台子羽居楚。子夏居西河。子贡终於齐。如田子方。段干木。吴起。禽滑釐之属。皆受业於子夏之伦。为王者师。是时独魏文侯好学。後陵迟以至于始皇。天下并争於战国。懦术既绌焉。然齐鲁之间。学者独不废也。於威。宣之际。孟子。荀卿之列。咸遵夫子之业而润色之。以学显於当世。\\
\hspace*{2em}及至秦之季世。焚诗书。坑术士。六艺从此缺焉。陈涉之王也。而鲁诸儒持孔氏之礼器往归陈王。於是孔甲为陈涉博士。卒与涉俱死。陈涉起匹夫。驱瓦合適戍。旬月以王楚。不满半岁竟灭亡。其事至微浅。然而缙绅先生之徒负孔子礼器往委质为臣者。何也。以秦焚其业。积怨而发愤于陈王也。\\
\hspace*{2em}及高皇帝诛项籍。举兵围鲁。鲁中诸儒尚讲诵习礼乐。弦歌之音不绝。岂非圣人之遗化。好礼乐之国哉。故孔子在陈。曰﹁归与归与。吾党之小子狂简。斐然成章。不知所以裁之﹂。夫齐鲁之间於文学。自古以来。其天性也。故汉兴。然後诸儒始得脩其经艺。讲习大射乡饮之礼。叔孙通作汉礼仪。因为太常。诸生弟子共定者。咸为选首。於是喟然叹兴於学。然尚有干戈。平定四海。亦未暇遑庠序之事也。孝惠。吕后时。公卿皆武力有功之臣。孝文时颇徵用。然孝文帝本好刑名之言。及至孝景。不任儒者。而窦太后又好黄老之术。故诸博士具官待问。未有进者。\\
\hspace*{2em}及今上即位。赵绾。王臧之属明儒学。而上亦乡之。於是招方正贤良文学之士。自是之後。言诗於鲁则申培公。於齐则辕固生。於燕则韩太傅。言尚书自济南伏生。言礼自鲁高堂生。言易自菑川田生。言春秋於齐鲁自胡毋生。於赵自董仲舒。及窦太后崩。武安侯田蚡为丞相。绌黄老。刑名百家之言。延文学儒者数百人。而公孙弘以春秋白衣为天子三公。封以平津侯。天下之学士靡然乡风矣。\\
\hspace*{2em}公孙弘为学官。悼道之郁滞。乃请曰。﹁丞相御史言。制曰﹃盖闻导民以礼。风之以乐。婚姻者。居屋之大伦也。今礼废乐崩。朕甚愍焉。故详延天下方正博闻之士。咸登诸朝。其令礼官劝学。讲议洽闻兴礼。以为天下先。太常议。与博士弟子。崇乡里之化。以广贤材焉﹄。谨与太常臧。博士平等议曰。闻三代之道。乡里有教。夏曰校。殷曰序。周曰庠。其劝善也。显之朝廷。其惩恶也。加之刑罚。故教化之行也。建首善自京师始。由内及外。今陛下昭至德。开大明。配天地。本人伦。劝学脩礼。崇化厉贤。以风四方。太平之原也。古者政教未洽。不备其礼。请因旧官而兴焉。为博士官置弟子五十人。复其身。太常择民年十八已上。仪状端正者。补博士弟子。郡国县道邑有好文学。敬长上。肃政教。顺乡里。出入不悖所闻者。令相长丞上属所二千石。二千石谨察可者。当与计偕。诣太常。得受业如弟子。一岁皆辄试。能通一艺以上。补文学掌故缺。其高弟可以为郎中者。太常籍奏。即有秀才异等。辄以名闻。其不事学若下材及不能通一艺。辄罢之。而请诸不称者罚。臣谨案诏书律令下者。明天人分际。通古今之义。文章尔雅。训辞深厚。恩施甚美。小吏浅闻。不能究宣。无以明布谕下。治礼次治掌故。以文学礼义为官。迁留滞。请选择其秩比二百石以上。及吏百石通一艺以上。补左右内史。大行卒史。比百石已下。补郡太守卒史。皆各二人。边郡一人。先用诵多者。若不足。乃择掌故补中二千石属。文学掌故补郡属。备员。请著功令。佗如律令。﹂制曰。﹁可。﹂自此以来。则公卿大夫士吏斌斌多文学之士矣。\\
\hspace*{2em}申公者。鲁人也。高祖过鲁。申公以弟子从师入见高祖于鲁南宫。吕太后时。申公游学长安。与刘郢同师。已而郢为楚王。令申公傅其太子戊。戊不好学。疾申公。及王郢卒。戊立为楚王。胥靡申公。申公耻之。归鲁。退居家教。终身不出门。复谢绝宾客。独王命召之乃往。弟子自远方至受业者百馀人。申公独以诗经为训以教。无传。疑者则阙不传。\\
\hspace*{2em}兰陵王臧既受诗。以事孝景帝为太子少傅。免去。今上初即位。臧乃上书宿卫上。累迁。一岁中为郎中令。及代赵绾亦尝受诗申公。绾为御史大夫。绾。臧请天子。欲立明堂以朝诸侯。不能就其事。乃言师申公。於是天子使使束帛加璧安车驷马迎申公。弟子二人乘轺传从。至。见天子。天子问治乱之事。申公时已八十馀。老。对曰。﹁为治者不在多言。顾力行何如耳。﹂是时天子方好文词。见申公对。默然。然已招致。则以为太中大夫。舍鲁邸。议明堂事。太皇窦太后好老子言。不说儒术。得赵绾。王臧之过以让上。上因废明堂事。尽下赵绾。王臧吏。後皆自杀。申公亦疾免以归。数年卒。\\
\hspace*{2em}弟子为博士者十馀人。孔安国至临淮太守。周霸至胶西内史。夏宽至城阳内史。砀鲁赐至东海太守。兰陵缪生至长沙内史。徐偃为胶西中尉。邹人阙门庆忌为胶东内史。其治官民皆有廉节。称其好学。学官弟子行虽不备。而至於大夫。郎中。掌故以百数。言诗虽殊。多本於申公。\\
\hspace*{2em}清河王太傅辕固生者。齐人也。以治诗。孝景时为博士。与黄生争论景帝前。黄生曰。﹁汤武非受命。乃弑也。﹂辕固生曰。﹁不然。夫桀纣虐乱。天下之心皆归汤武。汤武与天下之心而诛桀纣。桀纣之民不为之使而归汤武。汤武不得已而立。非受命为何。﹂黄生曰。﹁冠虽敝。必加於首。履虽新。必关於足。何者。上下之分也。今桀纣虽失道。然君上也。汤武虽圣。臣下也。夫主有失行。臣下不能正言匡过以尊天子。反因过而诛之。代立践南面。非弑而何也。﹂辕固生曰。﹁必若所云。是高帝代秦即天子之位。非邪。﹂於是景帝曰。﹁食肉不食马肝。不为不知味。言学者无言汤武受命。不为愚。﹂遂罢。是後学者莫敢明受命放杀者。\\
\hspace*{2em}窦太后好老子书。召辕固生问老子书。固曰。﹁此是家人言耳。﹂太后怒曰。﹁安得司空城旦书乎。﹂乃使固入圈刺豕。景帝知太后怒而固直言无罪。乃假固利兵。下圈刺豕。正中其心。一刺。豕应手而倒。太后默然。无以复罪。罢之。居顷之。景帝以固为廉直。拜为清河王太傅。久之。病免。\\
\hspace*{2em}今上初即位。复以贤良徵固。诸谀儒多疾毁固。曰﹁固老﹂。罢归之。时固已九十馀矣。固之徵也。薛人公孙弘亦徵。侧目而视固。固曰。﹁公孙子。务正学以言。无曲学以阿世。﹂自是之後。齐言诗皆本辕固生也。诸齐人以诗显贵。皆固之弟子也。\\
\hspace*{2em}韩生者。燕人也。孝文帝时为博士。景帝时为常山王太傅。韩生推诗之意而为内外传数万言。其语颇与齐鲁间殊。然其归一也。淮南贲生受之。自是之後。而燕赵间言诗者由韩生。韩生孙商为今上博士。\\
\hspace*{2em}伏生者。济南人也。故为秦博士。孝文帝时。欲求能治尚书者。天下无有。乃闻伏生能治。欲召之。是时伏生年九十馀。老。不能行。於是乃诏太常使掌故朝错往受之。秦时焚书。伏生壁藏之。其後兵大起。流亡。汉定。伏生求其书。亡数十篇。独得二十九篇。即以教于齐鲁之间。学者由是颇能言尚书。诸山东大师无不涉尚书以教矣。\\
\hspace*{2em}伏生教济南张生及欧阳生。欧阳生教千乘兒宽。兒宽既通尚书。以文学应郡举。诣博士受业。受业孔安国。兒宽贫无资用。常为弟子都养。及时时间行佣赁。以给衣食。行常带经。止息则诵习之。以试第次。补廷尉史。是时张汤方乡学。以为奏谳掾。以古法议决疑大狱。而爱幸宽。宽为人温良。有廉智。自持。而善著书。书奏。敏於文。口不能发明也。汤以为长者。数称誉之。及汤为御史大夫。以兒宽为掾。荐之天子。天子见问。说之。张汤死後六年。兒宽位至御史大夫。九年而以官卒。宽在三公位。以和良承意从容得久。然无有所匡谏。於官。官属易之。不为尽力。张生亦为博士。而伏生孙以治尚书徵。不能明也。\\
\hspace*{2em}自此之後。鲁周霸。孔安国。雒阳贾嘉。颇能言尚书事。孔氏有古文尚书。而安国以今文读之。因以起其家。逸书得十馀篇。盖尚书滋多於是矣。\\
\hspace*{2em}诸学者多言礼。而鲁高堂生最本。礼固自孔子时而其经不具。及至秦焚书。书散亡益多。於今独有士礼。高堂生能言之。\\
\hspace*{2em}而鲁徐生善为容。孝文帝时。徐生以容为礼官大夫。传子至孙延。徐襄。襄。其天姿善为容。不能通礼经。延颇能。未善也。襄以容为汉礼官大夫。至广陵内史。延及徐氏弟子公户满意。桓生。单次。皆尝为汉礼官大夫。而瑕丘萧奋以礼为淮阳太守。是後能言礼为容者。由徐氏焉。\\
\hspace*{2em}自鲁商瞿受易孔子。孔子卒。商瞿传易。六世至齐人田何。字子庄。而汉兴。田何传东武人王同子仲。子仲传菑川人杨何。何以易。元光元年徵。官至中大夫。齐人即墨成以易至城阳相。广川人孟但以易为太子门大夫。鲁人周霸。莒人衡胡。临菑人主父偃。皆以易至二千石。然要言易者本於杨何之家。\\
\hspace*{2em}董仲舒。广川人也。以治春秋。孝景时为博士。下帷讲诵。弟子传以久次相受业。或莫见其面。盖三年董仲舒不观於舍园。其精如此。进退容止。非礼不行。学士皆师尊之。今上即位。为江都相。以春秋灾异之变推阴阳所以错行。故求雨闭诸阳。纵诸阴。其止雨反是。行之一国。未尝不得所欲。中废为中大夫。居舍。著灾异之记。是时辽东高庙灾。主父偃疾之。取其书奏之天子。天子召诸生示其书。有刺讥。董仲舒弟子吕步舒不知其师书。以为下愚。於是下董仲舒吏。当死。诏赦之。於是董仲舒竟不敢复言灾异。\\
\hspace*{2em}董仲舒为人廉直。是时方外攘四夷。公孙弘治春秋不如董仲舒。而弘希世用事。位至公卿。董仲舒以弘为从谀。弘疾之。乃言上曰。﹁独董仲舒可使相缪西王。﹂胶西王素闻董仲舒有行。亦善待之。董仲舒恐久获罪。疾免居家。至卒。终不治产业。以脩学著书为事。故汉兴至于五世之间。唯董仲舒名为明於春秋。其传公羊氏也。\\
\hspace*{2em}胡毋生。齐人也。孝景时为博士。以老归教授。齐之言春秋者多受胡毋生。公孙弘亦颇受焉。\\
\hspace*{2em}瑕丘江生为穀梁春秋。自公孙弘得用。尝集比其义。卒用董仲舒。\\
\hspace*{2em}仲舒弟子遂者。兰陵褚大。广川殷忠。温吕步舒。褚大至梁相。步舒至长史。持节使决淮南狱。於诸侯擅专断。不报。以春秋之义正之。天子皆以为是。弟子通者。至於命大夫。为郎。谒者。掌故者以百数。而董仲舒子及孙皆以学至大官。\\

% ===== 索隐述赞（司马贞《史记索隐》）· 统一样式 =====
% 题头「索隐述赞曰」五字并入反白框；正文每四句一列，定宽盒 \hfil 居列中；
% \fontsize{24}{41} 列距对齐全书界行（121+41k）；\clearpage 使述赞独占一页。
\clearpage
\shuzanhead{索隐述赞曰}\linebreak
{\shuzanbody\fontsize{24}{41}\selectfont
\hbox to \hsize{\hfil 孔氏之衰。经书绪乱。言诸六学。始自炎汉。\hfil}\linebreak
\hbox to \hsize{\hfil 著令立官。四方鸧腕。曲台坏壁。书礼之冠。\hfil}\linebreak
\hbox to \hsize{\hfil 传易言诗。云蒸雾散。兴化致理。鸿猷克赞。\hfil}%
\par}
\end{document}
